\documentclass[10pt]{article}
\usepackage[T1]{fontenc}
\usepackage{lmodern,amsmath,amssymb,amsthm,microtype,booktabs}
\usepackage[margin=0.9in]{geometry}
\usepackage[colorlinks=true,urlcolor=blue,citecolor=blue]{hyperref}
\newtheorem{theorem}{Theorem}
\newtheorem{lemma}{Lemma}
\newtheorem{corollary}{Corollary}
\newcommand{\Om}{\Omega}
\newcommand{\om}{\omega}
\title{Prefix-Free Divisor Representations with Prescribed Prime Factorization}
\author{Qichao Wang\\{\small Hebei University of Technology}\\[0.3em]
Daoyu Dong\\{\small University of Electronic Science and Technology of China}}
\date{October 26, 2026}
\hypersetup{pdftitle={Prefix-Free Divisor Representations with Prescribed Prime Factorization},pdfauthor={Qichao Wang and Daoyu Dong}}
\begin{document}
\maketitle
\begin{abstract}
For every integer base $b\ge2$ and every finite list of positive prime
exponents, infinitely many integers with precisely that prime signature
have prefix-free base-$b$ representations of all their divisors greater
than one. This proves the existence conjectures recorded in OEIS A088862
and A090335. If the exponents sum to $r\ge1$, a witness is smaller than
$b^{r^3+32r}$. The construction separates divisor lengths and leading
values by choosing all prime factors in one short multiplicative interval.
\end{abstract}

\section{The two existence questions}
A finite collection of words is prefix-free if no word is a proper prefix
of another. Integer representations have no leading zeros. OEIS A088862
asks for the smallest integer with a prescribed total number $\Om(N)$ of
prime factors and prefix-free binary divisor representations; A090335
uses the number $\om(N)$ of distinct prime factors instead
\cite{a88862,a90335}. Both entries record the conjecture that a witness
exists for every prescribed count.

We use $r$ for the actual factor count. The displayed lists and synthesized
b-files currently begin with $1$, of factor count zero, while their
published first offset is $1$; their names, formulas and examples use the
factor count itself. Our statements include every count $r\ge0$ and thus
apply to either indexing convention. No list value is changed here.

\begin{lemma}[Separation interval]\label{lem:separate}
Fix $b\ge2$, $r\ge1$ and an integer $k\ge2$. Suppose that every prime factor,
with multiplicity, of an integer $N$ belongs to
\begin{equation}\label{eq:interval}
 I_{b,r,k}=\bigl(b^{k-1/r},\ b^{k-1/(r+1)}\bigr),
 \qquad \Om(N)=r.
\end{equation}
Then the base-$b$ representations of the divisors of $N$ greater than one
are prefix-free.
\end{lemma}
\begin{proof}
A divisor $d$ with $i=\Om(d)\in\{1,\ldots,r\}$ satisfies
\begin{equation}\label{eq:layer}
 b^{ik-i/r}<d<b^{ik-i/(r+1)}.
\end{equation}
Since $i/r\le1$, this places $d$ strictly between $b^{ik-1}$ and $b^{ik}$,
so its representation has exactly $ik$ digits. Two distinct divisors in
the same layer have the same length and cannot be proper prefixes.

Let $u,v$ belong to layers $i<j$. The integer represented by the first
$ik$ digits of $v$ is $\lfloor v/b^{(j-i)k}\rfloor$. But
\[
 \frac{v}{b^{(j-i)k}}<b^{ik-j/(r+1)}
 <b^{ik-i/r}<u,
\]
because $r(j-i)>i$ for $1\le i<j\le r$. The leading digits of $v$ therefore
cannot equal $u$.
\end{proof}

\newpage
\begin{lemma}[Enough primes, with an explicit scale]\label{lem:primes}
For every $b\ge2$, $r\ge1$ and integer $k\ge r^2+32$, the interval
$I_{b,r,k}$ contains at least $r$ distinct primes.
\end{lemma}
\begin{proof}
Dusart's Proposition 6.8 \cite{dusart} states that for $x>396738$ there
is a prime in $(x,x(1+1/(25\log^2x))]$. Put
\[
 L=b^{k-1/r},\quad U=b^{k-1/(r+1)},\quad
 \varepsilon=\frac1{25\log^2L}.
\]
Here $L\ge2^{32}>396738$. Starting at $L$, apply this result successively
at the prime just found. The $r$-th prime is at most
$L(1+\varepsilon)^r<L\exp(r\varepsilon)$. Also
\[
 25(k-1/r)^2\log^3b>\frac{25}{8}r^4>r^2(r+1),
\]
where $\log2>1/2$ and $k-1/r>r^2$ were used. Consequently
\[
 r\varepsilon<\frac{\log b}{r(r+1)}=\log(U/L).
\]
All $r$ primes lie strictly below $U$, as required.
\end{proof}

\begin{theorem}[Arbitrary prime signature]\label{thm:main}
Let $e_1,\ldots,e_t$ be positive integers and $r=e_1+\cdots+e_t$.
For each base $b\ge2$ there are infinitely many choices of distinct primes
$p_1,\ldots,p_t$ for which $N=\prod_{i=1}^t p_i^{e_i}$ has prefix-free
base-$b$ representations of its divisors greater than one. One such $N$
satisfies $N<b^{r^3+32r}$.
\end{theorem}
\begin{proof}
Take any $k\ge r^2+32$. Lemma~\ref{lem:primes} supplies at least $r\ge t$
distinct primes in $I_{b,r,k}$. Their product with the specified exponents
satisfies Lemma~\ref{lem:separate}. For $k=r^2+32$,
$N<b^{rk-r/(r+1)}<b^{r^3+32r}$. As $k$ tends to infinity the chosen primes,
and hence $N$, are unbounded.
\end{proof}

\begin{corollary}
For every prescribed factor count $r\ge0$, both existence questions in
A088862 and A090335 have affirmative answers. For $r\ge1$ each admits
infinitely many witnesses. In the first problem witnesses can be prime
powers $p^r$; in the second they can be squarefree products of $r$ primes.
\end{corollary}
\begin{proof}
For $r=0$ use $N=1$, whose set of divisors greater than one is empty.
For $r\ge1$ apply Theorem~\ref{thm:main} respectively to the exponent lists
$(r)$ and $(1,\ldots,1)$. A nonempty set of positive integers has a least
element.
\end{proof}

\section{Exact arithmetic and verification}
Interval membership can be checked without logarithms or floating-point
roots: an integer $p$ lies in $I_{b,r,k}$ exactly when
\[
 p^r>b^{kr-1},\qquad p^{r+1}<b^{k(r+1)-1}.
\]
The accompanying standard-library Python program constructs small
witnesses using these two integer inequalities. It independently
enumerates divisors and compares their digit strings, tests repeated
and mixed exponent lists, and verifies the initial published values.
The explicit bound uses Dusart's theorem; the finite checks audit the
construction and the transcription of the OEIS definitions.

\begin{thebibliography}{9}
\bibitem{a88862} R. Zumkeller, OEIS A088862,
\url{https://oeis.org/A088862}, accessed October 26, 2026.
\bibitem{a90335} R. Zumkeller, OEIS A090335,
\url{https://oeis.org/A090335}, accessed October 26, 2026.
\bibitem{dusart} P. Dusart, \emph{Estimates of some functions over primes
without R.H.}, arXiv:1002.0442v1 (2010), Proposition 6.8,
\url{https://arxiv.org/abs/1002.0442v1}.
\end{thebibliography}
\end{document}
